% chapters/abstract.tex
\chapter*{Abstract}
\addcontentsline{toc}{chapter}{Abstract}

% TARGET: 1 page (approximately 250-350 words)
%
% STRUCTURE:
% - Context (1-2 sentences)
% - Problem (2-3 sentences)
% - Methods (2-3 sentences)
% - Results (3-4 sentences)
% - Conclusions (2-3 sentences)

% ============================================================
% DRAFT TEMPLATE - FILL IN AFTER RESULTS ARE KNOWN
% ============================================================

% CONTEXT:
% Motor Imagery Brain-Computer Interfaces (MI-BCI) offer promise for 
% stroke rehabilitation by enabling patients to practice motor actions 
% through neural signals alone. The Apple Catcher system, a gamified 
% MI-BCI paradigm, previously achieved 78\% cross-subject classification 
% accuracy using EEGNet, a compact deep learning architecture.

% PROBLEM:
% However, this high accuracy raises concerns about validity. The system 
% uses a visual paradigm where an apple falls toward the target hand 
% before the formal cue, potentially introducing eye movement artifacts 
% or visual processing confounds. The finding that the pre-cue window 
% dramatically outperformed the post-cue window is counterintuitive and 
% unexplained. Without understanding what drives classification, the 
% system's clinical applicability remains uncertain.

% METHODS:
% This thesis conducts a systematic interpretability analysis of EEGNet 
% trained on the Apple Catcher dataset (40 subjects, 200 trials each). 
% We employ three complementary approaches: (1) neurophysiological 
% characterization of the dataset using time-frequency analysis for 
% event-related desynchronization and slow potentials for movement-related 
% cortical potentials; (2) visualization of learned temporal and spatial 
% filters; and (3) attribution analysis using Integrated Gradients to 
% identify decision-relevant features. Results are evaluated against 
% competing hypotheses: artifacts, classical ERD, motor preparation 
% (MRCP), and visual processing.

% RESULTS:
% [TO BE COMPLETED AFTER ANALYSIS]
% Our analysis reveals that [primary finding]. Temporal filters 
% predominantly target [frequency bands], while spatial filters focus 
% on [brain regions]. Attribution analysis shows that classification 
% decisions are primarily driven by [features] localized to [regions] 
% during the [time period]. Correlation analysis indicates [relationship 
% between ERD/accuracy].

% CONCLUSIONS:
% [TO BE COMPLETED AFTER ANALYSIS]
% These findings suggest that the Apple Catcher's classification 
% performance is [primarily valid / partially confounded / concerning]. 
% We provide specific recommendations for [deployment / further 
% development / paradigm modification] and demonstrate a methodology 
% for validating deep learning models in visual MI-BCI applications.

% ============================================================
% KEYWORDS
% ============================================================
\vspace{1cm}
\noindent\textbf{Keywords:} Brain-Computer Interface, Motor Imagery, 
Deep Learning, EEGNet, Interpretability, Integrated Gradients, 
Event-Related Desynchronization, Transfer Learning
